\documentclass[10pt,a4paper]{amsart}
\usepackage[margin=19mm]{geometry}
\usepackage{amsmath,amssymb,mathtools}
\usepackage[scr=rsfs,cal=euler]{mathalfa}
\usepackage{xfrac}
\usepackage[unicode,hidelinks,bookmarks=false]{hyperref}
\newcommand{\ie}{i.e.\ }
\newcommand{\cf}{cf.\ }
\newcommand{\resp}{respectively\ }
\newcommand{\Subsection}{\subsection}
\providecommand{\nobreakdash}{\nobreak\hbox{-}}
\setlength{\emergencystretch}{2em}
\multlinegap0pt
\usepackage{amssymb}
\usepackage{mathtools}
\usepackage{dsfont}
\usepackage{xcolor}
\usepackage{braket}
\usepackage[shortlabels]{enumitem}
\newtheorem{proposition}{Proposition}
\newcommand{\lset}{\left\{ }
\newcommand{\rset}{\right\} }
\newcommand{\mb}{\mathbb}
\newcommand{\mc}{\mathcal}
\DeclareMathOperator{\tr}{tr}
\title{How Far do Lindbladians Go?}
\author{Jihong Cai, Advith Govindarajan, Marius Junge}
\date{Source: arXiv:2504.04883v2 (CC BY 4.0)}
\begin{document}
\maketitle

\noindent\emph{Source and licence.} How Far do Lindbladians Go?. Version arXiv:2504.04883v2 (CC BY 4.0), distributed by arXiv under the Creative Commons Attribution 4.0 International licence, http://creativecommons.org/licenses/by/4.0/, as recorded in the arXiv licence metadata for this exact version and shown on the abstract page https://arxiv.org/abs/2504.04883v2. Full licence notice: \texttt{licenses/arxiv-2504.04883-CC-BY-4.0.txt}.\par\smallskip

\noindent\textit{Editorial scope.} Counted unit: the article's proposition nonconv (source0.tex lines 580--582). The article's complete original proof of it is delivered with it (source0.tex lines 583--607, one \texttt{proof} environment of 25 lines). No mathematics is written, repaired or re-derived: the statement and the proof are the article's own lines, in the article's own notation, and no retained line is substituted or re-flowed.  No further source-local context is needed: every cross-reference of every retained line resolves inside the two delivered spans.  Nothing was omitted from any retained span, so the package carries no omission record.  No retained line contains a citation, so the wrapper carries no bibliography and no bibliography entry.  No retained line contains a figure environment, an image-inclusion command or drawing code, so no image is delivered and the package declares no asset dependency.  Editorial formatting only, in addition: an AMS \texttt{amsart} preamble that restates the article's own declarations and macros used by the retained lines, namely the environments \texttt{proposition} on separate counters, together with the standard packages those lines need.  The retained environments print in this document as Proposition 1 in order of appearance, whereas the article prints the same items with the numbering of its own full document.  The complete native source of the article as distributed in the arXiv e-print for this exact version is bundled unchanged as \texttt{sources/176087/source0.tex}.
\begin{proposition}  \label{nonconv} For non-invertible state $\rho \in \partial \mc D(\mc{H})$
    $$T^+_\rho \mc D(\mc{H})\supsetneq \lset x\in\mb B(H): x=x^*, \tr x=0, \rho+\varepsilon x\geq 0\rset.$$
\end{proposition}

\begin{proof}
    The inclusion is immediate since if $\rho + \varepsilon x \geq 0$ for some $\varepsilon > 0$, then the curve $\rho + \varepsilon x$ lies in the state space $\mc D(\mc{H})$ for small $\varepsilon$, so $x \in T^+_\rho \mc D(\mc{H})$.
    
    To show that equality does not hold, we construct an example where a Hermitian traceless operator $x$ lies in the tangent cone, but $\rho + \varepsilon x \not\geq 0$ for any $\varepsilon > 0$.
    Let
    $$\rho = \begin{pmatrix} 1 & 0 \\ 0 & 0 \end{pmatrix}, \quad
    x = \begin{pmatrix} 0 & 1 \\ 1 & 0 \end{pmatrix}.$$
    Then
    $$\rho + \varepsilon x = \begin{pmatrix} 1 & \varepsilon \\ \varepsilon & 0 \end{pmatrix}$$
    has determinant $\det(\rho + \varepsilon x) = -\varepsilon^2 < 0$, so it is not positive semidefinite for any $\varepsilon > 0$. Hence, $x$ is not in the set on the right-hand side.
    
    However, define
    $$b = \begin{pmatrix} -2 & 0 \\ 0 & 2 \end{pmatrix},$$
    and consider the second-order perturbation:
    $$\rho_\varepsilon = \rho + \varepsilon x + \varepsilon^2 b = 
    \begin{pmatrix}
    1 - 2\varepsilon^2 & \varepsilon \\
    \varepsilon & 2\varepsilon^2
    \end{pmatrix}.$$
    The determinant of this matrix is
    $$\det(\rho_\varepsilon) = (1 - 2\varepsilon^2)(2\varepsilon^2) - \varepsilon^2 = 2\varepsilon^2 - 4\varepsilon^4 - \varepsilon^2 = \varepsilon^2(1 - 4\varepsilon^2).$$
    For small $\varepsilon > 0$, this is positive, and the matrix is clearly Hermitian with trace $1$. Hence, $\rho_\varepsilon \in \mc D(\mc{H})$ for small $\varepsilon$, and the curve $\rho_\varepsilon$ lies in the state space. Its derivative at $\varepsilon = 0$ is $x$, so $x \in T^+_\rho \mc D(\mc{H})$.
    
    Therefore, $x \in T^+_\rho \mc D(\mc{H})$, but $x$ is not in the set $\lset  x : \rho + \varepsilon x \geq 0 \rset $, showing that the inclusion is strict.
\end{proof}

\end{document}
